\documentclass[a4paper]{article}

% Basic packages
% Español (México) con XeLaTeX: fontspec para fuentes Unicode, polyglossia
% para la división silábica y los nombres del documento en la variante mexicana.
\usepackage{fontspec}
\usepackage{polyglossia}
\setdefaultlanguage[variant=mexican]{spanish}
% Los nombres chinos van como «pinyin (original)»: xeCJK da a los caracteres
% Han su propia fuente; sin ella XeLaTeX los omite con un aviso.
% FandolSong no cubre algunos caracteres de nombres (U+73FA); WenQuanYi Zen Hei sí.
\usepackage[AutoFallBack=true]{xeCJK}
\setCJKmainfont{FandolSong-Regular.otf}
\setCJKfallbackfamilyfont{\CJKrmdefault}{WenQuanYi Zen Hei}
% polyglossia no traduce los nombres de listings.
\AtBeginDocument{\providecommand{\lstlistingname}{}\renewcommand{\lstlistingname}{Listado}%
  \providecommand{\lstlistlistingname}{}\renewcommand{\lstlistlistingname}{Índice de listados}}
\usepackage{amsmath, amssymb}
\usepackage{graphicx}
\usepackage[margin=2.5cm]{geometry}
\usepackage[most]{tcolorbox}
\usepackage{etoolbox}
\usepackage{listings}
\usepackage{hyperref}
\usepackage{booktabs}
\usepackage{subcaption}
\usepackage{float}
\usepackage{tikz}
\IfFileExists{pgfplots.sty}{
    \usepackage{pgfplots}
    \pgfplotsset{compat=1.18}
}{}

% Highlight boxes
\newtcolorbox{knowledgebox}[1]{
    enhanced, colback=blue!5!white, colframe=blue!75!black, colbacktitle=blue!75!black,
    coltitle=white, fonttitle=\bfseries, title={#1}, attach boxed title to top left={yshift=-2mm, xshift=2mm},
    boxrule=1pt, sharp corners
}

\newtcolorbox{importantbox}[1]{
    enhanced, colback=yellow!10!white, colframe=yellow!80!black, colbacktitle=yellow!80!black,
    coltitle=black, fonttitle=\bfseries, title={#1}, sharp corners
}

\newtcolorbox{warningbox}[1]{
    enhanced, colback=red!5!white, colframe=red!75!black, colbacktitle=red!75!black,
    coltitle=white, fonttitle=\bfseries, title={#1}, sharp corners
}

% Listings
\lstset{
    language=Python,
    basicstyle=\ttfamily\small,
    keywordstyle=\color{blue},
    stringstyle=\color{red!60!black},
    commentstyle=\color{green!60!black},
    breaklines=true,
    frame=single,
    numbers=left,
    numberstyle=\tiny\color{gray},
    captionpos=b,
    extendedchars=true
}

% Front-page metadata
\newcommand{\notetitle}{MIT 6.S191 Lecture 6: Language Models and New Frontiers}
\newcommand{\noteauthors}{Elaboradas a partir de la clase de Alexander Amini}
\newcommand{\notedate}{Primavera de 2025}
\newcommand{\videochannel}{Alexander Amini (MIT)}
\newcommand{\videopublishdate}{2025-04-07}
\newcommand{\videoduration}{47 minutos}
\newcommand{\videourl}{https://www.youtube.com/watch?v=HLKo4fJx_7k}
\newcommand{\videocoverpath}{cover.jpg}
\newcommand{\repourl}{https://github.com/hqhq1025/ai-course-notes}

% Footer with repo info
\usepackage{fancyhdr}
\pagestyle{fancy}
\fancyhf{}
\fancyfoot[C]{\thepage}
\fancyfoot[R]{\footnotesize\color{gray}\href{\repourl}{AI Course Notes}}
\renewcommand{\headrulewidth}{0pt}
\renewcommand{\footrulewidth}{0pt}

\begin{document}

\begin{titlepage}
\centering
{\Large Notas del curso\par}
\vspace{1.2cm}
{\huge\bfseries \notetitle\par}
\vspace{0.8cm}
{\large \noteauthors\par}
\vspace{0.3cm}
{\large \notedate\par}
\vspace{1.2cm}

\ifdefempty{\videocoverpath}{
{\small Al generar las notas, indique la ruta local de la portada del video.\par}
}{
\includegraphics[width=0.82\textwidth,height=0.45\textheight,keepaspectratio]{\videocoverpath}\par
}

\vfill
\begin{tcolorbox}[width=0.9\textwidth, colback=black!2!white, colframe=black!60, sharp corners]
\textbf{Autor o canal del video}: \videochannel\par
\textbf{Fecha de publicación}: \videopublishdate\par
\textbf{Duración del video}: \videoduration\par
\ifdefempty{\videourl}{}{
\textbf{Enlace del video}: \href{\videourl}{\nolinkurl{\videourl}}\par
}\par
\textbf{Página del proyecto}: \href{\repourl}{\nolinkurl{\repourl}}\par
\end{tcolorbox}
\end{titlepage}

\tableofcontents
\newpage

%% --- Inicio del contenido --- %%

\section{Introducción: repaso del curso y propósito de esta clase}

Esta clase es la última de las seis clases fundamentales de la serie MIT 6.S191 Introduction to Deep Learning. En las cinco clases anteriores, el curso presentó de manera sistemática la metodología central del aprendizaje profundo: desde las redes neuronales completamente conectadas básicas, el modelado de secuencias (RNN / Transformer), las redes neuronales convolucionales (CNN), hasta los modelos generativos (VAE / GAN) y el aprendizaje por refuerzo profundo. Esta clase, apoyándose en estos fundamentos, se concentrará en dos temas centrales:

\begin{enumerate}
    \item \textbf{las limitaciones del aprendizaje profundo} (Limitations): el dilema de la generalización, los ataques adversarios, la representación de la incertidumbre, entre otros;
    \item \textbf{las nuevas fronteras} (New Frontiers): los avances más recientes representados por Diffusion Models y Large Language Models (LLMs).
\end{enumerate}

\begin{figure}[H]
\centering
\includegraphics[width=0.85\textwidth]{slides-images/slide-01.jpg}
\caption{Tema de esta clase: Deep Learning Limitations and New Frontiers\protect\footnotemark}
\end{figure}
\footnotetext{Slides, página 1.}

Como señaló Amini en la clase, el aprendizaje profundo ya ha traído avances revolucionarios en numerosos campos, como la conducción autónoma, las imágenes médicas, el aprendizaje por refuerzo, el modelado generativo, el procesamiento de lenguaje natural, las finanzas y la seguridad. Pero como profesionales de la tecnología, no basta con entender el poder de estos algoritmos; también debemos reconocer con claridad sus limitaciones, para poder desplegar la IA de manera responsable en el mundo real.

\begin{figure}[H]
\centering
\includegraphics[width=0.85\textwidth]{slides-images/slide-10.jpg}
\caption{Auge del aprendizaje profundo en diversos campos\protect\footnotemark}
\end{figure}
\footnotetext{Slides, página 10.}

\subsection{Resumen de la sección}

Esta clase, como cierre de las clases fundamentales, conecta lo anterior con lo que sigue: repasa la esencia del aprendizaje profundo como aproximador de funciones, señala sus limitaciones y presenta las nuevas fronteras representadas por Diffusion Models y LLM.

%% =====================================================
\section{La esencia de las redes neuronales: aproximadores de funciones}

\subsection{Universal Approximation Theorem}

Para entender el poder y las limitaciones del aprendizaje profundo, un punto de partida muy importante es el \textbf{Universal Approximation Theorem} (teorema de aproximación universal) propuesto por Hornik y colegas en 1989:

\begin{importantbox}{Universal Approximation Theorem}
Una red neuronal prealimentada de una sola capa oculta, si la cantidad de unidades ocultas es suficientemente grande, puede aproximar cualquier función continua con una precisión arbitraria.
\end{importantbox}

$$
f(x) \approx \hat{f}(x; W, b) = \sigma(W_2 \cdot \sigma(W_1 x + b_1) + b_2)
$$

Donde:
\begin{itemize}
    \item $f(x)$: función continua objetivo
    \item $\hat{f}(x; W, b)$: salida de la red neuronal con una sola capa oculta
    \item $\sigma$: función de activación no lineal
    \item $W_1, W_2, b_1, b_2$: parámetros de pesos y sesgos aprendibles
\end{itemize}

\begin{figure}[H]
\centering
\includegraphics[width=0.85\textwidth]{slides-images/slide-12.jpg}
\caption{Diagrama esquemático del Universal Approximation Theorem\protect\footnotemark}
\end{figure}
\footnotetext{Diapositiva 12. Cita a Hornik et al., Neural Networks, 1989.}

\subsection{Las limitaciones del teorema}

Aunque este teorema es muy potente en el plano teórico, existen dos caveats clave:

\begin{figure}[H]
\centering
\includegraphics[width=0.85\textwidth]{slides-images/slide-13.jpg}
\caption{Dos advertencias importantes del Universal Approximation Theorem\protect\footnotemark}
\end{figure}
\footnotetext{Diapositiva 13.}

\begin{warningbox}{Las dos advertencias principales del teorema de aproximación universal}
\begin{enumerate}
    \item \textbf{La cantidad de unidades ocultas puede ser inviablemente grande}: el teorema garantiza la existencia de dicha red, pero no da una cota superior para la cantidad de neuronas necesarias: para funciones complejas, podría requerirse una cantidad astronómica de parámetros.
    \item \textbf{No garantiza la capacidad de generalización}: el teorema solo garantiza la capacidad de aproximación, no que el modelo aprendido se desempeñe bien con datos no vistos. Ajustarse a los datos de entrenamiento no equivale a comprender las regularidades subyacentes de los datos.
\end{enumerate}
\end{warningbox}

Además, el teorema no ofrece una metodología para encontrar los pesos óptimos. Si bien el descenso de gradiente es una estrategia de optimización eficaz, en superficies de pérdida no convexas de alta dimensión, la optimización en sí misma es un problema sumamente no trivial.

\subsection{Perspectiva histórica del desarrollo de la IA}

\begin{figure}[H]
\centering
\includegraphics[width=0.85\textwidth]{slides-images/slide-14.jpg}
\caption{La trayectoria histórica del ``Hype'' de la IA\protect\footnotemark}
\end{figure}
\footnotetext{Diapositiva 14.}

El desarrollo de la IA atravesó un hype cycle típico: desde el nacimiento en la conferencia de Dartmouth en 1956, pasando por el primer invierno de la IA de 1974 a 1980, luego el segundo invierno de 1987 a 1993, hasta llegar finalmente a un crecimiento explosivo en la era del aprendizaje profundo. Comprender esta historia nos ayuda a mantener un juicio sereno sobre el auge actual de la IA: hay que abrazar el avance tecnológico, pero también estar alerta ante el exceso de exaltación.

\subsection{Resumen de la sección}

En esencia, las redes neuronales son aproximadores de funciones. El teorema de aproximación universal garantiza la capacidad expresiva en el plano teórico, pero en la aplicación práctica, el tamaño de la capa oculta, la dificultad de la optimización y la capacidad de generalización son desafíos reales que hay que enfrentar. Este marco de comprensión sienta las bases para la discusión, en la siguiente sección, de las limitaciones concretas del aprendizaje profundo.

%% =====================================================
\section{Limitaciones del aprendizaje profundo}

\subsection{Revisitando la generalización: el experimento de etiquetas aleatorias}

Amini presenta un experimento clásico, del artículo de Zhang et al., publicado en 2017 en ICLR: \textit{Understanding Deep Neural Networks Requires Rethinking Generalization}.

\begin{knowledgebox}{Diseño del experimento}
Los investigadores tomaron las imágenes del conjunto de datos ImageNet junto con sus etiquetas y, para cada imagen de forma independiente, lanzaron un dado de $k$ caras ($k$ es el número total de categorías), sustituyendo la etiqueta original por el resultado aleatorio. La consecuencia de esto es que dos imágenes de la misma categoría pueden terminar asignadas a etiquetas completamente distintas, y la asociación semántica entre la etiqueta y el contenido de la imagen se pierde por completo.
\end{knowledgebox}

\begin{figure}[H]
\centering
\includegraphics[width=0.85\textwidth]{slides-images/slide-19.jpg}
\caption{Experimento de etiquetas aleatorias: las etiquetas de las imágenes se mezclan por completo\protect\footnotemark}
\end{figure}
\footnotetext{Diapositiva 19. Cita a Zhang+ ICLR 2017.}

\subsubsection{Resultados del experimento}

Los investigadores entrenaron redes neuronales profundas con distintos grados de aleatorización de las etiquetas, observando el desempeño en el conjunto de entrenamiento y en el conjunto de prueba:

\begin{figure}[H]
\centering
\includegraphics[width=0.85\textwidth]{slides-images/slide-22.jpg}
\caption{Hallazgo clave: incluso con etiquetas completamente aleatorias, la precisión de entrenamiento puede llegar al 100\%\protect\footnotemark}
\end{figure}
\footnotetext{Diapositiva 22. La línea punteada roja marca que la precisión de entrenamiento se mantiene siempre cerca del 100\%.}

\begin{importantbox}{Hallazgo central: la sorprendente capacidad de ajuste de las redes profundas}
\begin{itemize}
    \item \textbf{Precisión de prueba}: conforme aumenta el grado de aleatorización de las etiquetas, la precisión en el conjunto de prueba disminuye continuamente (algo intuitivo).
    \item \textbf{Precisión de entrenamiento}: sin importar qué tan aleatorias sean las etiquetas, la red profunda siempre puede alcanzar una precisión cercana al 100\% en el conjunto de entrenamiento.
\end{itemize}
Esto significa que las redes profundas actuales tienen capacidad suficiente para ``memorizar'' todo el conjunto de entrenamiento, incluso cuando en los datos no existe ningún patrón que se pueda aprender.
\end{importantbox}

\subsection{La doble filo del aproximador de funciones}

A partir del experimento anterior, Amini profundiza en este fenómeno desde la perspectiva del ajuste de funciones:

\begin{figure}[H]
\centering
\includegraphics[width=0.85\textwidth]{slides-images/slide-25.jpg}
\caption{La red neuronal como aproximador de funciones: un desempeño excelente en las regiones con datos de entrenamiento\protect\footnotemark}
\end{figure}
\footnotetext{Diapositiva 25.}

La red neuronal puede ajustar muy bien los valores de la función cerca de los puntos de datos de entrenamiento. Dado un nuevo punto de datos (en morado), si cae dentro del rango de la distribución de los datos de entrenamiento, la red puede dar una predicción razonable. Pero el problema central es el siguiente:

\begin{figure}[H]
\centering
\includegraphics[width=0.85\textwidth]{slides-images/slide-28.jpg}
\caption{El comportamiento fuera de la distribución de los datos es impredecible\protect\footnotemark}
\end{figure}
\footnotetext{Diapositiva 28.}

\begin{warningbox}{El problema out-of-distribution}
Cuando los datos de entrada quedan fuera de la distribución de entrenamiento (out-of-distribution, OOD), el comportamiento de la red neuronal es impredecible. En las regiones sin datos de entrenamiento, la función ajustada puede producir valores de salida arbitrarios, y la red no te ``avisa'' que no está segura. Esto plantea una pregunta central: \textbf{¿cómo sabemos qué es lo que el modelo no sabe?}
\end{warningbox}

\subsection{El aprendizaje profundo no es alquimia}

\begin{figure}[H]
\centering
\includegraphics[width=0.85\textwidth]{slides-images/slide-29.jpg}
\caption{Deep Learning $\neq$ Alchemy\protect\footnotemark}
\end{figure}
\footnotetext{Diapositiva 29. Cita a U. Muller, 6.S191 2018.}

Amini enfatiza: el aprendizaje profundo \textbf{no} es una ``caja negra mágica'' en la que puedes meter cualquier cosa y obtener una salida perfecta. El principio clásico de ``garbage in, garbage out'' (basura entra, basura sale) se aplica por completo. Antes de desplegar una solución de IA hay que considerar dos preguntas centrales:
\begin{enumerate}
    \item ¿Esta tarea realmente es adecuada para resolverse con aprendizaje profundo?
    \item ¿Puedes recopilar y depurar datos de alta calidad para esta tarea?
\end{enumerate}

\subsection{Modo de falla uno: sesgo en los datos}

\begin{figure}[H]
\centering
\begin{subfigure}[b]{0.48\textwidth}
    \includegraphics[width=\textwidth]{slides-images/slide-30.jpg}
    \caption{En las imágenes de entrenamiento, una gran cantidad de perros aparecen con la lengua afuera}
\end{subfigure}
\hfill
\begin{subfigure}[b]{0.48\textwidth}
    \includegraphics[width=\textwidth]{slides-images/slide-31.jpg}
    \caption{Zonas rosas que aparecen cuando la CNN colorea imágenes en blanco y negro}
\end{subfigure}
\caption{Falla causada por el sesgo en los datos: el modelo de coloreado de imágenes\protect\footnotemark}
\end{figure}
\footnotetext{Diapositivas 30--31. Cita a P. Isola 6.869.}

El curso presenta un ejemplo elocuente: entrenar una CNN para convertir fotografías en blanco y negro a color. Cuando en la fotografía de salida de un perro aparecía una mancha rosa en la zona de la mandíbula, la razón era que, en el conjunto de entrenamiento, una gran cantidad de fotografías de perros los mostraban con la lengua rosa afuera; el modelo aprendió la asociación estadística de que ``cerca de la mandíbula del perro debería haber rosa'', y no una verdadera comprensión del significado del color.

\subsection{Modo de falla dos: incertidumbre en escenarios críticos para la seguridad}

\begin{figure}[H]
\centering
\includegraphics[width=0.85\textwidth]{slides-images/slide-32.jpg}
\caption{Caso de accidente fatal del Autopilot de Tesla\protect\footnotemark}
\end{figure}
\footnotetext{Diapositiva 32. Cita a ABC News.}

Amini menciona un accidente fatal real de conducción autónoma: un Tesla, en el mismo tramo de carretera, se desvió repetidamente hacia la barrera divisoria hasta que finalmente terminó en una colisión. La investigación reveló que la barrera divisoria de ese tramo se había construido después de la recopilación de los datos de entrenamiento: el sistema de conducción autónoma se encontró con un escenario \textbf{fuera de la distribución de entrenamiento}.

\begin{figure}[H]
\centering
\includegraphics[width=0.85\textwidth]{slides-images/slide-33.jpg}
\caption{La incertidumbre es crucial en las aplicaciones críticas para la seguridad\protect\footnotemark}
\end{figure}
\footnotetext{Diapositiva 33.}

\begin{importantbox}{La incertidumbre en el aprendizaje profundo}
En los siguientes escenarios, detectar y cuantificar la incertidumbre de manera confiable es especialmente importante:
\begin{itemize}
    \item \textbf{Aplicaciones críticas para la seguridad}: conducción autónoma, diagnóstico médico, reconocimiento facial
    \item \textbf{Problemas de calidad de los datos}: desbalance de clases (imbalance), ruido en los datos (data noise)
    \item \textbf{Detección fuera de distribución}: cuando los datos de entrada difieren de manera significativa de la distribución de entrenamiento, el modelo debería ``saber que no sabe''
\end{itemize}
\end{importantbox}

\subsection{Modo de falla tres: ataques adversarios}

Los ataques adversarios (Adversarial Attacks) son un problema de seguridad clásico y muy importante en el aprendizaje profundo.

\begin{figure}[H]
\centering
\includegraphics[width=0.85\textwidth]{slides-images/slide-35.jpg}
\caption{Ejemplo de ataque adversario: Temple (97\%) $\rightarrow$ Ostrich (98\%)\protect\footnotemark}
\end{figure}
\footnotetext{Diapositiva 35.}

\subsubsection{El fundamento matemático de los ataques adversarios}

Recordemos la fórmula central del descenso de gradiente:

$$
W \leftarrow W - \eta \frac{\partial J(W, x, y)}{\partial W}
$$

donde $W$ son los pesos, $\eta$ es la tasa de aprendizaje, $J$ es la función de pérdida y $(x, y)$ son la entrada y la etiqueta fijas. Durante el entrenamiento, \textbf{fijamos la entrada $x$ y la etiqueta $y$}, y optimizamos los pesos $W$ para minimizar la pérdida.

\begin{figure}[H]
\centering
\includegraphics[width=0.85\textwidth]{slides-images/slide-40.jpg}
\caption{Fórmula para generar la imagen adversaria\protect\footnotemark}
\end{figure}
\footnotetext{Diapositiva 40.}

El ataque adversario invierte por completo este proceso:

$$
x \leftarrow x + \eta \frac{\partial J(W, x, y)}{\partial x}
$$

\begin{importantbox}{Ataque adversario vs. entrenamiento normal}
\begin{itemize}
    \item \textbf{Entrenamiento normal}: se fija la entrada $(x, y)$ y se optimizan los pesos $W$ para \textbf{minimizar} la pérdida
    \item \textbf{Ataque adversario}: se fijan los pesos $W$ y la etiqueta $y$, y se modifica la entrada $x$ para \textbf{maximizar} la pérdida
\end{itemize}
La diferencia clave: el objeto respecto al cual se calcula el gradiente pasa de $\partial W$ a $\partial x$, y la dirección de la optimización pasa de un signo menos a un signo más.
\end{importantbox}

\subsubsection{Ejemplos adversarios en el mundo físico}

Los ataques adversarios no solo existen en el espacio digital, también pueden realizarse en el mundo físico. Amini presenta un ejemplo adversario físico fabricado por estudiantes del MIT con impresión 3D: una tortuga impresa en 3D con una textura cuidadosamente diseñada que, desde diversos ángulos y bajo distintas condiciones de iluminación, siempre es clasificada por el clasificador como ``rifle'' en lugar de ``tortuga''.

\begin{figure}[H]
\centering
\includegraphics[width=0.85\textwidth]{slides-images/slide-42.jpg}
\caption{Ejemplo adversario físico impreso en 3D: una tortuga identificada como rifle\protect\footnotemark}
\end{figure}
\footnotetext{Diapositiva 42. Cita a Athalye+ ICML 2018.}

\begin{warningbox}{La amenaza real de los ataques adversarios}
La existencia de ejemplos adversarios implica que, en escenarios críticos para la seguridad (como el reconocimiento de señales de tránsito en la conducción autónoma o el reconocimiento facial en sistemas de seguridad), una perturbación mínima y cuidadosamente diseñada puede provocar que el sistema emita un juicio completamente erróneo. Esto no es solo un problema académico, sino un riesgo de seguridad que debe considerarse necesariamente al desplegar sistemas de IA en la práctica.
\end{warningbox}

\subsection{Síntesis de las limitaciones}

\begin{figure}[H]
\centering
\includegraphics[width=0.85\textwidth]{slides-images/slide-46.jpg}
\caption{Lista completa de las limitaciones de las redes neuronales\protect\footnotemark}
\end{figure}
\footnotetext{Diapositiva 46.}

El curso resume las principales limitaciones de las redes neuronales:

\begin{table}[H]
\centering
\begin{tabular}{ll}
\toprule
\textbf{Limitación} & \textbf{Descripción} \\
\midrule
Hambrienta de datos (Data hungry) & Por lo general requiere millones de muestras de entrenamiento \\
Intensiva en cómputo (Computationally intensive) & Tanto el entrenamiento como el despliegue requieren GPU \\
Vulnerable a ataques adversarios (Adversarial examples) & Fácil de engañar con perturbaciones cuidadosamente diseñadas \\
Sesgo algorítmico (Algorithmic bias) & Puede amplificar los sesgos sociales presentes en los datos \\
Mala representación de la incertidumbre (Uncertainty) & Es difícil saber cuándo el modelo ``no sabe'' \\
Caja negra inexplicable (Black boxes) & Difícil de entender y de confiar en ella \\
Requiere conocimiento experto (Expert knowledge) & El diseño de la arquitectura y el ajuste de hiperparámetros dependen de la experiencia \\
\textbf{Dificultad para codificar estructura} (Encode structure) & Difícil incorporar conocimiento previo al aprendizaje \\
\textbf{Dificultad de extrapolación} (Extrapolation) & Difícil ir más allá del rango de la distribución de los datos de entrenamiento \\
\bottomrule
\end{tabular}
\caption{Principales limitaciones de las redes neuronales}
\end{table}

\begin{knowledgebox}{Las limitaciones engendran nuevas fronteras}
Amini señala: como técnicos y científicos, las limitaciones representan precisamente oportunidades de desarrollo. Los Diffusion Models y los LLM que se presentan en la segunda mitad del curso son justamente una respuesta a estos dos retos centrales: la ``dificultad de extrapolación'' y la ``dificultad para codificar estructura''; al modelar la propia distribución de los datos, se espera que los modelos base (Foundation Models) logren una capacidad de generalización más sólida.
\end{knowledgebox}

\subsection{Resumen de la sección}

Las limitaciones del aprendizaje profundo abarcan múltiples aspectos: una capacidad de generalización insuficiente (el experimento de etiquetas aleatorias reveló una sorprendente capacidad de sobreajuste), errores sistemáticos causados por el sesgo en los datos, retos de incertidumbre en escenarios críticos para la seguridad, y la amenaza que representan los ataques adversarios para la robustez del modelo. Comprender estas limitaciones es un requisito para usar la tecnología de IA de manera responsable, y también es un motor que impulsa el desarrollo de la próxima generación de algoritmos.

%% =====================================================
\section{Nueva frontera I: IA generativa y modelos de difusión}

\subsection{Trayectoria del desarrollo del modelado generativo}

En la Lecture 4, el curso presentó dos modelos generativos clásicos: el VAE (Variational Autoencoder) y la GAN (Generative Adversarial Network). Aunque estos dos tipos de modelos lograron, de manera pionera, la capacidad de generar nuevas muestras a partir de datos, enfrentan varias limitaciones clave:

\begin{figure}[H]
\centering
\includegraphics[width=0.85\textwidth]{slides-images/slide-48.jpg}
\caption{Limitaciones y retos del VAE/GAN\protect\footnotemark}
\end{figure}
\footnotetext{Diapositiva 48.}

\begin{knowledgebox}{Principales problemas del VAE y la GAN}
\textbf{Limitaciones (Limitations)}:
\begin{itemize}
    \item \textbf{Mode collapse}: el modelo tiende a generar únicamente muestras ``promedio'', perdiendo diversidad
    \item \textbf{Dificultad para generar fuera de la distribución}: es difícil generar muestras novedosas fuera de la distribución de entrenamiento
    \item \textbf{Dificultad de entrenamiento}: en particular, el proceso de entrenamiento adversario de la GAN es sumamente inestable
\end{itemize}
\textbf{Retos (Challenges)}: estabilidad, eficiencia, calidad de generación, novedad
\end{knowledgebox}

\subsection{Idea central de los Diffusion Models}

Los Diffusion Models proponen un paradigma de generación completamente distinto al del VAE/GAN:

\begin{figure}[H]
\centering
\includegraphics[width=0.85\textwidth]{slides-images/slide-50.jpg}
\caption{Generación ``one-shot'' del VAE/GAN vs. generación por eliminación iterativa de ruido de Diffusion\protect\footnotemark}
\end{figure}
\footnotetext{Diapositiva 50.}

\begin{importantbox}{Innovación clave de los Diffusion Models}
\begin{itemize}
    \item \textbf{VAE/GAN}: genera la muestra completa \textbf{de una sola vez} (one-shot) a partir de una variable latente de baja dimensión
    \item \textbf{Diffusion Models}: genera la muestra eliminando el ruido \textbf{de manera iterativa} (iteratively) y por pasos
\end{itemize}
Esta estrategia iterativa descompone una tarea de generación difícil en una serie de pasos simples de eliminación de ruido, lo que mejora notablemente la calidad de generación.
\end{importantbox}

\subsection{Proceso de adición de ruido hacia adelante (Forward Noising)}

El primer paso de un modelo de Diffusion es construir los datos de entrenamiento mediante un proceso de adición de ruido hacia adelante (Forward Noising Process).

\begin{figure}[H]
\centering
\includegraphics[width=0.85\textwidth]{slides-images/slide-51.jpg}
\caption{Proceso de difusión: adición de ruido hacia adelante y eliminación de ruido hacia atrás\protect\footnotemark}
\end{figure}
\footnotetext{Diapositiva 51. Cita a Sohl-Dickstein y otros, ICML 2015; Ho y otros, NeurIPS 2020.}

\begin{figure}[H]
\centering
\includegraphics[width=0.85\textwidth]{slides-images/slide-53.jpg}
\caption{Proceso de adición de ruido hacia adelante: adición progresiva de ruido\protect\footnotemark}
\end{figure}
\footnotetext{Diapositiva 53.}

Pasos clave del proceso hacia adelante:
\begin{enumerate}
    \item Se parte de una imagen de entrenamiento $x_0$
    \item Se toma una muestra de un patrón de ruido aleatorio
    \item Se aumenta gradualmente la proporción de ruido según un programa (schedule) predefinido:
    \begin{itemize}
        \item $T=0$: 100\% imagen, 0\% ruido
        \item $T=1$: 75\% imagen, 25\% ruido
        \item $T=2$: 50\% imagen, 50\% ruido
        \item $T=3$: 25\% imagen, 75\% ruido
        \item $T=4$: 0\% imagen, 100\% ruido (ruido puramente aleatorio)
    \end{itemize}
\end{enumerate}

\begin{knowledgebox}{El proceso hacia adelante no requiere entrenamiento}
El proceso de adición de ruido hacia adelante es completamente determinista (dado un schedule de ruido) y no involucra ningún parámetro entrenable. Su único propósito es generar pares de entrenamiento para el proceso de eliminación de ruido hacia atrás: cada par $(x_t, x_{t-1})$ constituye una muestra de entrenamiento.
\end{knowledgebox}

\subsection{Proceso de eliminación de ruido hacia atrás (Reverse Denoising)}

La tarea central del entrenamiento es que la red neuronal aprenda el proceso inverso: dada la imagen con ruido en el instante $T$, estimar la imagen más limpia del instante $T-1$.

\begin{figure}[H]
\centering
\includegraphics[width=0.85\textwidth]{slides-images/slide-54.jpg}
\caption{Eliminación de ruido hacia atrás: aprender a estimar $T-1$ a partir de $T$\protect\footnotemark}
\end{figure}
\footnotetext{Diapositiva 54. Cita a Ho y otros, NeurIPS 2020.}

\begin{importantbox}{Objetivo de entrenamiento de Diffusion}
Dada la imagen con ruido $x_T$ en el instante $T$, se entrena una red neuronal $f_\theta$ para predecir la imagen $x_{T-1}$ del instante $T-1$:
$$
\hat{x}_{T-1} = f_\theta(x_T, T)
$$
La pérdida de entrenamiento es la diferencia entre la imagen predicha y la imagen real (habitualmente el MSE, o el MSE del ruido predicho). Esta red comparte parámetros en todos los pasos de tiempo, y el paso de tiempo $T$ se usa como entrada de condicionamiento.
\end{importantbox}

\subsection{Muestreo para generar imágenes nuevas}

Una vez terminado el entrenamiento, el proceso de generar una imagen nueva consiste en partir de ruido puro y aplicar repetidamente la red de eliminación de ruido:

\begin{figure}[H]
\centering
\includegraphics[width=0.85\textwidth]{slides-images/slide-61.jpg}
\caption{Proceso de muestreo: de ruido puro a una imagen nítida, paso a paso\protect\footnotemark}
\end{figure}
\footnotetext{Diapositiva 61.}

Proceso de muestreo: $x_T \xrightarrow{f_\theta} x_{T-1} \xrightarrow{f_\theta} x_{T-2} \xrightarrow{f_\theta} \cdots \xrightarrow{f_\theta} x_0$

Cada paso de eliminación de ruido solo necesita retirar ``una pequeña parte'' del ruido, lo cual es mucho más sencillo que generar la imagen completa de una sola vez; esta es precisamente la razón fundamental de la calidad superior de generación de los Diffusion Models.

\subsection{De la imagen al lenguaje natural: Text-to-Image}

Una aplicación importante de los Diffusion Models es su combinación con el procesamiento de lenguaje natural: la generación de imágenes a partir de texto (Text-to-Image Generation).

\begin{figure}[H]
\centering
\includegraphics[width=0.85\textwidth]{slides-images/slide-64.jpg}
\caption{Generación de imágenes a partir de descripciones en lenguaje natural\protect\footnotemark}
\end{figure}
\footnotetext{Diapositiva 64. Cita a Ramesh y otros, arXiv 2022.}

\begin{figure}[H]
\centering
\includegraphics[width=0.85\textwidth]{slides-images/slide-65.jpg}
\caption{Ejemplos diversos de generación Text-to-Image\protect\footnotemark}
\end{figure}
\footnotetext{Diapositiva 65. Cita a OpenAI, Ramesh y otros, arXiv 2022.}

Los sistemas Text-to-Image dominantes actuales (como DALL-E, Stable Diffusion, Midjourney) usan mayoritariamente un Diffusion Model como red troncal de generación de imágenes. El texto introducido por el usuario se transforma, mediante un codificador de lenguaje, en un vector de condicionamiento que guía el proceso de eliminación de ruido para que genere la imagen en la dirección que corresponde a la descripción textual.

\subsection{Más allá de la imagen: diseño molecular}

\begin{figure}[H]
\centering
\includegraphics[width=0.85\textwidth]{slides-images/slide-66.jpg}
\caption{Aplicación de los Diffusion Models en el diseño molecular y la generación de proteínas\protect\footnotemark}
\end{figure}
\footnotetext{Diapositiva 66.}

Las aplicaciones de los Diffusion Models van mucho más allá de la generación de imágenes. En los campos de la química y la biología, se emplea el mismo paradigma de ``del ruido a la estructura'' para:
\begin{itemize}
    \item \textbf{Generación de moléculas}: generar de manera progresiva estructuras de moléculas de fármacos a partir de ruido 3D
    \item \textbf{Diseño de proteínas}: generar nuevas estructuras de proteínas con funciones específicas
\end{itemize}

\begin{knowledgebox}{Generalidad de Diffusion}
El marco de ``adición de ruido $\rightarrow$ eliminación de ruido'' de los Diffusion Models tiene una generalidad muy alta: basta con que los datos puedan ser ``corrompidos'' y ``restaurados'' de manera progresiva para que este marco pueda aplicarse. Esto lo convierte en uno de los paradigmas de modelado generativo más activos en la actualidad, con aplicaciones que abarcan imágenes, video, audio, estructuras 3D, moléculas y muchos otros campos.
\end{knowledgebox}

\subsection{Resumen de la sección}

Los Diffusion Models, al descomponer la tarea de generación en múltiples pasos iterativos de eliminación de ruido, superan las limitaciones del VAE/GAN en cuanto a calidad de generación, estabilidad y diversidad. Sus ideas centrales incluyen: construir los datos de entrenamiento mediante la adición de ruido hacia adelante, aprender el proceso de generación mediante la eliminación de ruido hacia atrás y, finalmente, muestrear muestras de alta calidad a partir de ruido puro. Este paradigma se ha aplicado ampliamente en campos como Text-to-Image, la generación de video y el diseño molecular.

%% =====================================================
\section{Nueva frontera II: modelos de lenguaje grande (LLMs)}

\subsection{Definición y ubicación de los LLM}

\begin{figure}[H]
\centering
\includegraphics[width=0.85\textwidth]{slides-images/slide-69.jpg}
\caption{Ubicación de los LLM en el campo de la IA\protect\footnotemark}
\end{figure}
\footnotetext{Diapositiva 69.}

\begin{importantbox}{¿Qué es un modelo de lenguaje grande (LLM)?}
Los LLM son un concepto jerárquico dentro del campo de la IA:
\begin{itemize}
    \item \textbf{Artificial Intelligence}: todas las técnicas que hacen que una computadora simule el comportamiento humano
    \item \textbf{Deep Learning}: el uso de redes neuronales para extraer patrones de los datos
    \item \textbf{Large Language Models}: \textbf{redes neuronales de escala muy grande} entrenadas sobre \textbf{datos de texto de escala muy grande}
\end{itemize}
Palabras clave: ``very, very large neural networks trained on very, very large sets of text''.
\end{importantbox}

\subsection{Cómo funcionan los LLM: Next Token Prediction}

\begin{figure}[H]
\centering
\includegraphics[width=0.85\textwidth]{slides-images/slide-70.jpg}
\caption{Flujo de entrenamiento de un LLM tipo GPT\protect\footnotemark}
\end{figure}
\footnotetext{Diapositiva 70.}

El entrenamiento de un LLM puede resumirse en un objetivo sumamente sencillo:

$$
\text{Given a sequence of tokens, predict the next token.}
$$

Flujo concreto:
\begin{enumerate}
    \item \textbf{Conjunto de datos}: corpus de texto a gran escala como Common Crawl o WebText, divididos en tokens
    \item \textbf{Modelo}: por ejemplo GPT-3, con 175,000 millones de parámetros
    \item \textbf{Objetivo de entrenamiento}: dada una secuencia de tokens, predecir la distribución de probabilidad del siguiente token
    \item \textbf{Función de pérdida}: Cross-Entropy Loss, que mide la diferencia entre la distribución predicha y el siguiente token real
\end{enumerate}

\begin{figure}[H]
\centering
\includegraphics[width=0.85\textwidth]{slides-images/slide-71.jpg}
\caption{Flujo detallado de entrenamiento de Next Token Prediction\protect\footnotemark}
\end{figure}
\footnotetext{Diapositiva 71.}

\begin{knowledgebox}{De Next Token Prediction a la generación de texto}
Una vez concluido el entrenamiento, el LLM se usa mediante \textbf{generación autorregresiva} (autoregressive generation):
\begin{enumerate}
    \item El usuario introduce un prompt (por ejemplo, ``I'm giving a talk on AI at MIT. Can you outline it?'')
    \item El modelo predice la distribución de probabilidad del siguiente token y muestrea un token
    \item El nuevo token se agrega al final de la secuencia y se repite el proceso anterior
    \item Hasta generar un token de fin o alcanzar la longitud máxima
\end{enumerate}
Toda la ``inteligencia'' emergente se sustenta, en esencia, en esta tarea aparentemente simple de next token prediction.
\end{knowledgebox}

\subsection{Capacidades de los LLM}

\begin{figure}[H]
\centering
\includegraphics[width=0.85\textwidth]{slides-images/slide-73.jpg}
\caption{Capacidades actualmente confiables de los LLM\protect\footnotemark}
\end{figure}
\footnotetext{Diapositiva 73.}

GPT y otros LLM ya muestran un ``dominio'' (mastery over natural language) del lenguaje natural, con un desempeño confiable en los siguientes aspectos:
\begin{itemize}
    \item \textbf{Recuperación de conocimiento} (Knowledge Retrieval): responder preguntas factuales
    \item \textbf{Asistencia en escritura} (Writing Co-Pilot): redacción, pulido y traducción de textos
    \item \textbf{Asistencia en planeación} (Planning Co-Pilot): elaboración de planes y diseño de experimentos
\end{itemize}

\subsection{Limitaciones de los LLM}

\begin{figure}[H]
\centering
\includegraphics[width=0.85\textwidth]{slides-images/slide-74.jpg}
\caption{Las cuatro grandes limitaciones de los LLM\protect\footnotemark}
\end{figure}
\footnotetext{Diapositiva 74.}

\begin{warningbox}{Los cuatro retos clave de los LLM}
\begin{enumerate}
    \item \textbf{Robustez} (Robustness): son extremadamente sensibles a errores de ortografía o cambios de formato en la entrada (por ejemplo, ``Cn @uN66rN you translate ths from Spanish to English?'')
    \item \textbf{Alucinaciones} (Hallucinations): generan información incorrecta con total seguridad. ``GPT is a language model that...'' puede parecer perfectamente razonable, pero el contenido puede ser completamente inventado
    \item \textbf{Barandales y jailbreaks} (Guardrails and Jailbreaks): las restricciones de seguridad a nivel del modelo pueden eludirse con estrategias ingeniosas
    \item \textbf{Lógica y razonamiento numérico} (Logic and Numerics): tienen un desempeño débil en tareas que requieren razonamiento lógico riguroso y cálculo matemático
\end{enumerate}
\end{warningbox}

La raíz común de estos retos está en el \textbf{proceso de pensamiento de alto nivel} de los LLM: la robustez y la calibración de la confianza, la capacidad de planeación de largo plazo, y el razonamiento lógico y el descubrimiento científico son, todos ellos, fronteras actuales de investigación.

\subsection{Capacidades emergentes y leyes de escalamiento}

\begin{figure}[H]
\centering
\includegraphics[width=0.85\textwidth]{slides-images/slide-75.jpg}
\caption{Capacidades que emergen al crecer la escala del modelo\protect\footnotemark}
\end{figure}
\footnotetext{Diapositiva 75. Cita a Wei y otros, TMLR 2022.}

\begin{importantbox}{Capacidades emergentes (Emergent Abilities)}
Se dice que una capacidad es \textbf{emergente} (emergent) si no existe en modelos pequeños pero aparece en modelos grandes. En concreto:
\begin{itemize}
    \item \textbf{Estructuración del lenguaje} (Structuring Language): surge de manera repentina cuando el modelo alcanza del orden de $\sim 10^{22}$ parámetros
    \item \textbf{Comprensión de la fonética} (Understanding Phonetics): un salto igualmente abrupto
    \item \textbf{Realización de operaciones aritméticas} (Performing Arithmetic): aparece solo a mayor escala
\end{itemize}
La aparición de estas capacidades presenta un rasgo de ``transición de fase'' (phase transition): no es una mejora gradual, sino que surge de golpe al cruzar cierto umbral de escala.
\end{importantbox}

\begin{figure}[H]
\centering
\includegraphics[width=0.85\textwidth]{slides-images/slide-76.jpg}
\caption{Árbol de capacidades emergentes: a partir de 8,000 millones de parámetros\protect\footnotemark}
\end{figure}
\footnotetext{Diapositiva 76. Cita a Google AI Blog.}

A medida que la escala del modelo crece a partir de 8,000 millones de parámetros, el LLM va adquiriendo, en orden, comprensión del lenguaje, operaciones aritméticas y capacidad de responder preguntas, como un árbol de capacidades que crece y florece sin cesar.

\subsection{Resumen de la sección}

Los LLM son, en esencia, redes neuronales de escala muy grande entrenadas sobre volúmenes masivos de datos de texto, que logran una capacidad lingüística sorprendente mediante el objetivo simple de next token prediction. Sus capacidades presentan un carácter emergente conforme crece la escala del modelo. Pero los LLM aún enfrentan retos importantes en robustez, alucinaciones, protecciones de seguridad y razonamiento lógico, que son también las direcciones de investigación más activas en la actualidad.

%% =====================================================
\section{Modelos fundacionales y la visión de la IA general}

\begin{figure}[H]
\centering
\includegraphics[width=0.85\textwidth]{slides-images/slide-77.jpg}
\caption{La idea poderosa que impulsan los Foundation Models\protect\footnotemark}
\end{figure}
\footnotetext{Diapositivas, página 77.}

Al final de la clase, Amini amplía la mirada hacia la visión más amplia de los Foundation Models:

\begin{importantbox}{La visión central de los Foundation Models}
\begin{itemize}
    \item \textbf{Sistema central de razonamiento}: ¿pueden los modelos fundacionales generativos ofrecer un motor de razonamiento unificado para la IA general?
    \item \textbf{IA que diseña IA}: ¿puede usarse la IA para mejorar y hacer evolucionar a la propia IA?
    \item \textbf{IA generativa multidominio}: de las imágenes a la biología, al lenguaje y más allá; con un enorme potencial, pero que exige cautela
    \item \textbf{La relación entre la inteligencia artificial y la inteligencia humana}: entender los vínculos y las diferencias entre ambas
\end{itemize}
\end{importantbox}

Estas preguntas no son solo desafíos técnicos, sino que también implican una reflexión científica y filosófica profunda. Como enfatiza Amini: el poder y la cautela van de la mano (power and caution); al mismo tiempo que impulsamos la frontera de la IA, debemos considerar con prudencia su impacto social.

%% =====================================================
\section{Síntesis y ampliación}

\subsection{Repaso de los puntos clave de esta clase}

Esta clase, como cierre de la serie de clases fundamentales de MIT 6.S191, cubrió tres temas centrales:

\begin{enumerate}
    \item \textbf{La naturaleza y las limitaciones del aprendizaje profundo}
    \begin{itemize}
        \item Las redes neuronales son aproximadores de funciones; el Universal Approximation Theorem garantiza la capacidad teórica, pero no la generalización
        \item El experimento con etiquetas aleatorias reveló la sorprendente capacidad de sobreajuste (memorización) de las redes profundas
        \item Tres modos de falla principales: sesgo de los datos, incertidumbre en escenarios críticos para la seguridad, ataques adversarios
        \item Nueve limitaciones principales: hambre de datos, uso intensivo de cómputo, fragilidad ante ataques adversarios, sesgo algorítmico, mala cuantificación de la incertidumbre, caja negra, necesidad de conocimiento experto, dificultad para codificar estructura, dificultad para extrapolar
    \end{itemize}

    \item \textbf{Nueva frontera I -- Diffusion Models}
    \begin{itemize}
        \item Superan el mode collapse y la inestabilidad de entrenamiento de las VAE/GAN
        \item Construyen pares de entrenamiento agregando ruido hacia adelante, y generan mediante eliminación de ruido hacia atrás
        \item La estrategia de generación iterativa es el núcleo de su ventaja en calidad
        \item Aplicaciones amplias: texto a imagen, generación de video, diseño de moléculas y proteínas
    \end{itemize}

    \item \textbf{Nueva frontera II -- Large Language Models}
    \begin{itemize}
        \item Su naturaleza: redes neuronales a escala masiva más datos de texto a escala masiva
        \item Objetivo de entrenamiento central: Next Token Prediction
        \item Las capacidades emergentes surgen con un crecimiento en forma de «cambio de fase» a medida que aumenta el tamaño del modelo
        \item Retos clave: robustez, alucinaciones, protección de seguridad, razonamiento lógico
    \end{itemize}
\end{enumerate}

\subsection{Lecturas adicionales}

\begin{itemize}
    \item Hornik, K., Stinchcombe, M., \& White, H. (1989). Multilayer feedforward networks are universal approximators. \textit{Neural Networks}.
    \item Zhang, C., et al. (2017). Understanding deep learning requires rethinking generalization. \textit{ICLR 2017}.
    \item Goodfellow, I., Shlens, J., \& Szegedy, C. (2015). Explaining and harnessing adversarial examples. \textit{ICLR 2015}.
    \item Athalye, A., et al. (2018). Synthesizing robust adversarial examples. \textit{ICML 2018}.
    \item Ho, J., Jain, A., \& Abbeel, P. (2020). Denoising diffusion probabilistic models. \textit{NeurIPS 2020}.
    \item Sohl-Dickstein, J., et al. (2015). Deep unsupervised learning using nonequilibrium thermodynamics. \textit{ICML 2015}.
    \item Ramesh, A., et al. (2022). Hierarchical text-conditional image generation with CLIP latents. \textit{arXiv 2022}.
    \item Wei, J., et al. (2022). Emergent abilities of large language models. \textit{TMLR 2022}.
    \item Sitio oficial del curso: \url{http://introtodeeplearning.com}
    \item GitHub: \url{https://github.com/MITDeepLearning/introtodeeplearning/}
\end{itemize}

%% --- Fin del contenido --- %%

\end{document}
